\chapter{A Concrete Naming Certificate for $\SuslinLineUp$}
\label{ch:concrete-instances-suslin-line-namecert}
\origin{human}

Following Kelley and Engelking's treatment of ordered spaces and separability, the $\SuslinLineUp$ packet records the traditional Suslin-line boundary as finite order-topology NameCert data. It reads the total-order source of \autoref{ch:concrete-instances-totalorder-namecert}, the topology surface of \autoref{ch:concrete-instances-topology-namecert}, the ordered-real comparison of \autoref{ch:concrete-instances-linearly-ordered-real-namecert}, and the separability boundary of \autoref{ch:concrete-instances-separable-space-namecert}. The packet names a dense linear order without endpoints, the order-topology base read, a countable-chain-condition row, a nonseparability boundary row, transport, replay, provenance, and local naming; it does not construct a host Suslin line, decide Suslin's hypothesis, import a choice-selected dense set, or identify the line with the Real carrier.

\begin{definition}[Suslin line carrier]
\label{def:suslin-line-carrier}
A $\SuslinLineUp$ carrier is a finite $\BHist$ packet
$$
S=(O,T,C,B,H,K,P,N)
$$
with the following displayed rows:
\begin{enumerate}
\item $O$ is the total-order source row, naming a dense linear order without endpoints and exposing the comparison ledger inherited from \autoref{ch:concrete-instances-totalorder-namecert}.
\item $T$ is the topology row, recording the order-topology base and its local open-interval read through \autoref{ch:concrete-instances-topology-namecert}.
\item $C$ is the countable-chain-condition row, stating that pairwise-disjoint nonempty open intervals are admitted only through the displayed countability ledger.
\item $B$ is the separability-boundary row, recording that no countable dense subset is exported by the packet.
\item $H$ stores componentwise $\hsame$ transport for the order source, topology row, countable-chain-condition row, separability-boundary row, continuation, provenance, and local naming rows.
\item $K$ records displayed $\Cont$ replay from order and topology requests through $O,T,C,B$.
\item $P$ is the $\Pkg$ provenance over $\SuslinLineUp$, the total-order, topology, ordered-real, separability, $\BHist$, $\hsame$, $\Cont$, and $\NameCert$ rows.
\item $N$ is the local $\NameCert_{\SuslinLineUp}$ row.
\end{enumerate}
The classifier compares accepted carriers only through $O,T,C,B,H,K,P,N$. It has no coordinate for host set-theoretic existence, a forcing extension, a selected dense enumeration, quotient Real equality, or a proof or refutation of Suslin's hypothesis.
\end{definition}

\begin{theorem}[Suslin line NameCert obligations]
\label{thm:suslin-line-namecert-obligations}
For every accepted $\SuslinLineUp$ carrier $S=(O,T,C,B,H,K,P,N)$, the local $\NameCert_{\SuslinLineUp}$ surface consists exactly of carrier admission for the total-order source, order-topology base, countable-chain-condition row, separability-boundary row, transport, replay, provenance, and local naming rows; order visibility through $O$; topology visibility through $T$; countable-chain-condition visibility through $C$; nonseparability-boundary visibility through $B$; componentwise transport through $H$; displayed continuation through $K$; package provenance through $P$; and local naming through $N$.
\end{theorem}
\begin{proof}
Project the carrier of \autoref{def:suslin-line-carrier}. The accepted inventory is exactly $O,T,C,B,H,K,P,N$, so every consumer must read the order source, topology row, countable-chain-condition row, and separability-boundary row before the local name is available.

The sibling anchors supply only their displayed surfaces. The total-order and topology chapters provide comparison and open-interval data, the ordered-real chapter supplies a comparison reference rather than an identification with $\RealUp$, and the separable-space chapter provides the boundary against a displayed countable dense subset.

Rows $H,K,P,N$ preserve and replay that same finite route. Since no carrier coordinate names a host construction, forcing extension, selected dense enumeration, quotient Real equality, or set-theoretic decision principle, none of those endpoints can be exported by the NameCert obligations.
\end{proof}

\begin{theorem}[Suslin line separability boundary]
\label{thm:suslin-line-separability-boundary}
Every topology, total-order, or ordered-real consumer of an accepted $\SuslinLineUp$ carrier must read the separability-boundary row $B$ after the displayed order, topology, and countable-chain-condition rows $O,T,C$. The packet therefore records the Suslin-line boundary as a finite order-topology naming certificate: countable-chain-condition data may be consumed as a displayed row, but no countable dense subset, selected dense enumeration, quotient Real equality, host Suslin-line construction, or decision of Suslin's hypothesis is exported.
\end{theorem}
\begin{proof}
Start with \autoref{thm:suslin-line-namecert-obligations}. It requires the order row $O$, topology row $T$, and countable-chain-condition row $C$ before the separability-boundary row can be consumed.

Now compare with the separable-space anchor. The row $B$ records the absence of an exported countable dense subset in the same finite packet that carries the order topology and countable-chain-condition rows, so a consumer cannot replace the boundary by a selected dense enumeration.

Finally apply the structural rows of \autoref{def:suslin-line-carrier}. Transport, replay, provenance, and local naming preserve only the displayed order-topology packet, and they add no host construction, quotient Real equality, forcing decision, or ambient choice principle.
\end{proof}

\falsifiablePrediction{Every accepted $\SuslinLineUp$ topology-facing read factors through a total-order source row, topology row, countable-chain-condition row, separability-boundary row, hsame transport, Cont replay, Pkg provenance, and local naming.}

\independenceWitness{totalorder, topology, linearly-ordered-real, separable-space: $\SuslinLineUp$ is not bijective to any listed sibling because it jointly carries a dense linear-order source, order-topology base, countable-chain-condition row, nonseparability-boundary row, transport, replay, provenance, and local NameCert rows.}

\closureat{SuslinLineUp}{seedStr}

\begin{closurestatus}{\SuslinLineUp}
  \constructivestory{A $\SuslinLineUp$ packet is generated from a displayed total-order source, order-topology row, countable-chain-condition row, separability-boundary row, componentwise hsame transport, displayed Cont replay, Pkg provenance, and a local NameCert row.}
  \theoryclosure{\seedClosure}
  \scopeclosed{\autoref{def:suslin-line-carrier}, \autoref{thm:suslin-line-namecert-obligations}, and \autoref{thm:suslin-line-separability-boundary} seed the finite Suslin-line order-topology packet.}
  \formalstatus{\unformalizedV}
  \bridgestatus{none}
  \notclaimed{This chapter does not close Lean verification, bridge checking, Suslin's hypothesis, host construction of a Suslin line, forcing independence, a choice-selected dense subset, quotient Real equality, or ambient set-theoretic completeness.}
  \upgradepath{Obligation closure requires checked carrier admission, order-source discipline, topology-base exactness, countable-chain-condition visibility, separability-boundary nonescape, transport, replay, provenance, and local naming for BEDC.Derived.SuslinLineUp.}
\end{closurestatus}
